% =====================================================================
%  Bloque para pegar en \subsection{Local minima and barren plateau}
%  (línea ~619 de main.tex, justo antes de \section{Conclusions})
%
%  Figuras que hay que subir a la carpeta del manuscrito:
%     fig_bp_1a_l1.pdf   fig_bp_1a_l2.pdf
%     fig_bp_k6_l1.pdf   fig_bp_k6_l2.pdf   (estas dos sólo si se incluye K6)
%
%  Usa el mismo estilo de figura del manuscrito (subfigure + overpic), así
%  que las etiquetas a) y b) las pone LaTeX y NO vienen dentro del PDF.
% =====================================================================

\subsection{Local minima and barren plateau}

A recurring concern for variational algorithms is whether the cost landscape
can be navigated at all. Two distinct obstructions are usually invoked: barren
plateaus, where the gradient variance vanishes exponentially with system size
and the landscape becomes featureless~\cite{LaroccaNatRevPhys2025,
HolmesPRXQuantum2022}, and rough landscapes densely populated with local
minima, in which the optimizer stalls even though gradients remain sizable.
The distinction matters here because the enlarged local Hilbert space of qudits
has been argued to amplify barren plateaus relative to the qubit
case~\cite{FriedrichQMI2025}, so the trainability of a qudit ansatz cannot be
taken for granted.

Qudit-ADAPT addresses these obstructions through its warm-start construction
rather than through the expressivity of a fixed ansatz. To quantify this, we
isolate the contribution of the warm start while keeping everything else fixed.
At every ADAPT iteration $k$ we re-optimize the \emph{same} $k$-parameter
ansatz from three different initial points:
%
\begin{enumerate}
    \item the warm start actually used by the algorithm,
          $\boldsymbol{\theta}_k = (\boldsymbol{\theta}_{k-1}^{\,*}, 0)$, where
          $\boldsymbol{\theta}_{k-1}^{\,*}$ are the optimized parameters of the
          previous iteration;
    \item a cold restart at the origin, $\boldsymbol{\theta}_k = \boldsymbol{0}_k$,
          which resets the state to $\ket{\phi_g}$ at every step;
    \item $100$ independent random initializations,
          $\theta_j \sim \mathcal{U}(-\pi,\pi)$.
\end{enumerate}
%
The operator sequence is in all three cases the one selected by the gradient
criterion, so the only variable is the starting point of the optimization. For
the random initializations we record only the converged optimum, so that each
mark in Fig.~\ref{fig5} is a minimum in which an optimization actually
terminated, and each column reveals the local-minimum structure seen by an
ansatz of that size.

% ---------------------------------------------------------------------
\begin{figure}[H]
    \centering

    \begin{subfigure}
        \centering
        \begin{overpic}[width=0.99\linewidth]{fig_bp_1a_l1.pdf}
            \put(2,88){a)}
        \end{overpic}
    \end{subfigure}

    \vspace{0.4em}

    \begin{subfigure}
        \centering
        \begin{overpic}[width=0.99\linewidth]{fig_bp_1a_l2.pdf}
            \put(2,88){b)}
        \end{overpic}
    \end{subfigure}

    \caption{Relative error as a function of the number of variational
    parameters for instance 1, using operator pools obtained from the
    a) $l=1$ and b) $l=2$ truncations of the approximate AGP. Each horizontal
    mark is the converged optimum of one of $100$ independent random parameter
    initializations of the same ansatz, colored by its relative error. The blue
    curve is the warm-start strategy used by Qudit-ADAPT and the red curve a
    cold restart at $\boldsymbol{\theta}_k=\boldsymbol{0}_k$. The graph
    instance is shown as an inset.}
    \label{fig5}
\end{figure}
% ---------------------------------------------------------------------

Figure~\ref{fig5} shows the result for instance 1. The landscape is manifestly
rough: for $l=2$ the $100$ random initializations converge to $100$ distinct
minima spanning more than two orders of magnitude in relative error, from
$8\times10^{-3}$ to $1.3\times10^{-1}$. This is not a landscape in which an
uninformed optimizer succeeds by chance. Nonetheless the warm-start curve
descends monotonically below the entire cloud, reaching
$\epsilon_{\rm rel}=3.1\times10^{-4}$ against a median of $5.5\times10^{-2}$
over the random restarts, an improvement of more than two orders of magnitude
with respect to the typical outcome of an uninformed optimization of the very
same ansatz. For $l=1$ the algorithm converges at $k=21$ with
$\epsilon_{\rm rel}=6.6\times10^{-3}$, against a median of $1.4\times10^{-2}$
and a worst case of $1.5\times10^{-1}$ among $88$ distinct minima.

The cold restart isolates the mechanism. It uses the same operator sequence and
the same pool, and differs from the warm start only in discarding the
parameters of the previous iteration. For $l=2$ the two curves coincide up to
$k\simeq22$ and separate afterwards, with the cold restart saturating near
$3.3\times10^{-3}$ while the warm start continues to improve by an additional
order of magnitude. The advantage therefore does not originate in the
counterdiabatic pool alone, but in the combination of that pool with parameter
recycling.

The reason is structural. Because $\theta_{k+1}=0$ leaves the state unchanged,
the enlarged landscape satisfies
$E_{k+1}(\boldsymbol{\theta},0) = E_k(\boldsymbol{\theta})$, so the warm start
inherits the value of the previous optimum. At that point all derivatives with
respect to the inherited parameters vanish, while the derivative along the new
direction is precisely the gradient maximized by the selection criterion in
Eq.~(\ref{eq:adapt_gradient}), which is nonzero by construction unless the
convergence threshold has been met. A configuration that was a local minimum in
$k$ dimensions is therefore a saddle point in $k+1$ dimensions, with a
guaranteed descent direction. Rather than escaping local traps by climbing over
them, the algorithm bypasses them through the dimension it has just added,
which is the burrowing behaviour reported for the qubit
case~\cite{GrimsleynpjQuantumInf2023}. Two consequences follow: the optimized
energy is monotonically non-increasing in $k$, and the local minima of the
fixed-size landscape are never visited, because the optimization is never
initialized at an arbitrary point of that landscape.

We stress that this argument concerns the roughness of the landscape and not
the expressivity of the pool. Burrowing guarantees monotonic improvement and
escape from the minima of the previous ansatz, but it does not prevent the
algorithm from terminating at a suboptimal point when the gradients of the
entire pool become small, which is what limits the $l=1$ results. Consistently,
the gradient variance evaluated at random parameter points shows no exponential
flattening with ansatz depth over the range explored, indicating that what traps
the random restarts are local minima rather than a barren plateau.

% ---------------------------------------------------------------------
\begin{figure}[H]
    \centering

    \begin{subfigure}
        \centering
        \begin{overpic}[width=0.99\linewidth]{fig_bp_k6_l1.pdf}
            \put(2,88){a)}
        \end{overpic}
    \end{subfigure}

    \vspace{0.4em}

    \begin{subfigure}
        \centering
        \begin{overpic}[width=0.99\linewidth]{fig_bp_k6_l2.pdf}
            \put(2,88){b)}
        \end{overpic}
    \end{subfigure}

    \caption{Same analysis as in Fig.~\ref{fig5} for the complete graph $K_6$,
    using the a) $l=1$ and b) $l=2$ operator pools. The larger connectivity
    increases the pool size from $66$ to $2724$ operators, and with the $l=2$
    pool the algorithm converges to machine precision after $18$ parameters,
    where all three initialization strategies become indistinguishable.}
    \label{fig6}
\end{figure}
% ---------------------------------------------------------------------

The size of the advantage depends on how expressive the pool is relative to the
instance. Figure~\ref{fig6} repeats the analysis for the complete graph $K_6$,
whose larger connectivity increases the pool from $66$ operators at $l=1$ to
$2724$ at $l=2$.

For $l=1$ the landscape remains rough: the $100$ random initializations still
converge to $99$ distinct minima spanning two orders of magnitude, and the warm
start improves on their median by a factor of about four ($1.5\times10^{-4}$
against $5.4\times10^{-4}$). The comparison is nevertheless less clean than for
instance~1: $14$ of the $100$ random restarts terminate below the warm-start
value, and the cold restart itself ends slightly lower. This is not in conflict
with the burrowing argument, which guarantees monotonic improvement and the
existence of a descent direction at every step, but not global optimality of
the $k$-parameter landscape. Starting from
$(\boldsymbol{\theta}_{k-1}^{\,*},0)$ confines the optimization to the basin
containing the previous optimum, so a deeper basin located elsewhere is
unreachable without first increasing the energy, whereas a random
initialization may land in it directly. The relevant point is that this is a
low-probability event that costs a full optimization each time it is attempted:
recovering the best of $100$ restarts requires two orders of magnitude more
optimizations than the single one performed by the warm start, and identifying
which restart is best requires prior knowledge of $E_0$.

For $l=2$ the picture changes qualitatively. The algorithm reaches machine
precision after $18$ parameters and warm, cold and random initializations all
converge to the same solution, with the residual differences at the level of
numerical noise. Once the ansatz is sufficiently expressive the landscape no
longer presents competing minima and the initialization becomes irrelevant.
Taken together, Figs.~\ref{fig5} and~\ref{fig6} indicate that the benefit of the
warm start is largest precisely in the regime of practical interest, where the
pool is deliberately kept compact in order to limit circuit depth, and that it
becomes immaterial only when the operator pool is large enough that the
optimization problem is no longer difficult.
